泛函行列式的固有时表示把一圈紫外问题转化为短固有时的局域渐近展开。对欧氏正谱二次算符 \(\Delta\)，先选取逆长度减除尺度 \(\mu_\kappa\,[\mathrm{m^{-1}}]\)。谱定理把标量恒等式 \(\ln(\lambda/\mu_\kappa^2)\) 提升为
\begin{equation}
 \boxed{
 \Tr\ln\frac{\Delta}{\mu_\kappa^2}
 =-\int_0^\infty\frac{\dd s}{s}
 \Tr\!\left(\ee^{-s\Delta}-\ee^{-s\mu_\kappa^2}\right).}
 \label{eq:proper-time-functional-dp11}
\end{equation}
因为 \(\Delta\) 的本征值具有 \(\mathrm{m^{-2}}\) 量纲，固有时 \(s\) 具有 \(\mathrm{m^2}\) 量纲；大本征值对应 \(s\to0^+\)，所以紫外发散完全由热核的短时系数决定。若改用物理动量尺度，则 \(\mu_p=\hbar\mu_\kappa\)，不能把二者直接混入同一个无量纲比值。

在平直欧氏空间中，对 Laplace 型算符定义
\begin{equation}
 \boxed{
 \Delta=-\left(\nabla^2+E\right),
 \qquad
 \nabla_\mu=\partial_\mu+\omega_\mu,
 \qquad
 \Omega_{\mu\nu}=[\nabla_\mu,\nabla_\nu].}
 \label{eq:laplace-type-operator-dp11}
\end{equation}
其中 \(E\) 是纤维内部空间上的局域端同态，\(\omega_\mu\) 是联络，\(\Omega_{\mu\nu}\) 是其曲率。对角热核的短时展开为
\begin{equation}
 \boxed{
 K(x,x;s)
 =\frac1{(4\pi s)^{d/2}}
 \left[
 I+sE+s^2\left(
 \frac12E^2+\frac1{12}\Omega_{\mu\nu}\Omega^{\mu\nu}
 +\frac16\nabla^2E
 \right)+\order(s^3)
 \right].}
 \label{eq:heat-kernel-a4-flat-dp11}
\end{equation}
无边界体积分中 \(\nabla^2E\) 是全导数，因此四维对数紫外极点真正依赖的局域不变量是 \(E^2\) 与 \(\Omega_{\mu\nu}\Omega^{\mu\nu}\)。

在 \(d=4-2\epsilon\) 中，\eqnrefs{eq:heat-kernel-a4-flat-dp11} 的 \(s^2\) 项与 \(\dd s/s\) 组合成 \(\dd s\,s^{-1+\epsilon}\)，从而直接产生 \(1/\epsilon\) 极点。对实玻色场的一圈高斯行列式，发散部分为
\begin{equation}
 \boxed{
 \left.\frac12\Tr\ln\Delta\right|_{\rm div}
 =-\frac1{2(4\pi)^2\bar\epsilon}
 \int\dd^4 x\,
 \tr\!\left[
 \frac12E^2+\frac1{12}\Omega_{\mu\nu}\Omega^{\mu\nu}
 \right].}
 \label{eq:heat-kernel-divergence-master-dp11}
\end{equation}
反交换标量变量和 Dirac 费米变量会因统计性质改变泛函行列式的整体符号与简并度，但允许出现的局域算符基仍由同一短时系数决定。对非阿贝尔规范场，规范场本身与规范固定 Jacobian 产生的 Grassmann 辅助场分别贡献各自的有限维指标迹；这些辅助变量的出现来自规范轨道换元的 Jacobian，而不是额外的物理渐近自由度。
